\clearpage
\section{부분체와 대수적 폐포}
\label{sec:finite-field-subfields}

\begin{learninggoal}
고정된 대수적 폐포 안에서 유한체 사이의 포함관계를 분류한다. 교집합과 합성체를 계산하고, 유한체의 모든 대수적 확장에서 분리성과 근의 완전한 분해가 성립함을 증명한다.
\end{learninggoal}

앞으로 소수 $p$를 고정하고 모든 유한체 $\F_{p^n}$을 하나의 대수적 폐포 $\overline{\F_p}$ 안의 부분체로 본다. 유한체의 유일성 때문에 이렇게 배치한 뒤에는 원소 수가 같은 부분체가 실제로 하나뿐이다.

\begin{lemma}
한 체 $\Omega$ 안의 두 유한부분체 $E,F\subseteq\Omega$가 같은 수의 원소를 가지면
\[
  E=F
\]
이다.
\end{lemma}

\begin{proof}
$|E|=|F|=q$라 하자. 두 체의 원소는 모두 $X^q-X$의 근이다. 이 다항식은 중근을 갖지 않고 $\Omega$ 안에서 근을 많아야 $q$개 가지므로, $E$와 $F$는 같은 근 집합이다.
\end{proof}

\begin{theorem}[유한체의 포함관계]
양의 정수 $m,n$에 대해
\[
  \F_{p^m}\subseteq\F_{p^n}
  \quad\Longleftrightarrow\quad
  m\mid n.
\]
이 경우
\[
  [\F_{p^n}:\F_{p^m}]=\frac{n}{m}.
\]
\end{theorem}

\begin{proof}
포함이 성립하면 탑 법칙에 의해
\[
  n=[\F_{p^n}:\F_p]
  =[\F_{p^n}:\F_{p^m}]\,m
\]
이므로 $m\mid n$이다.

반대로 $n=rm$이라 하자. $a\in\F_{p^m}$이면 $a^{p^m}=a$이고, 이 등식을 $r$번 반복하면
\[
  a^{p^n}=a^{p^{rm}}=a.
\]
따라서 $a$는 $X^{p^n}-X$의 근이고 $\F_{p^n}$에 속한다. 차수 공식은 탑 법칙에서 따른다.
\end{proof}

\begin{corollary}[부분체의 분류]
$\F_{p^n}$의 부분체는 정확히
\[
  \F_{p^d}\qquad(d\mid n)
\]
들이다. 각 약수 $d$에 대해 원소가 $p^d$개인 부분체는 하나뿐이다.
\end{corollary}

\begin{proof}
부분체 $E\subseteq\F_{p^n}$의 표수는 $p$이고 $|E|=p^d$인 어떤 $d$가 존재한다. 탑 법칙으로 $d\mid n$이며, 고정된 대수적 폐포 안에서 원소가 $p^d$개인 부분체는 $\F_{p^d}$ 하나뿐이다. 역으로 $d\mid n$이면 앞의 정리에 의해 $\F_{p^d}\subseteq\F_{p^n}$이다.
\end{proof}

\begin{example}
$\F_{2^{12}}$의 부분체들은 $12$의 약수에 대응한다.
\[
  \F_2,\quad \F_{2^2},\quad \F_{2^3},\quad
  \F_{2^4},\quad \F_{2^6},\quad \F_{2^{12}}.
\]
반면 $\F_{2^5}$는 $5\nmid12$이므로 $\F_{2^{12}}$의 부분체가 아니다.
\end{example}

\subsection{교집합과 합성체}

\begin{corollary}
고정된 $\overline{\F_p}$ 안에서
\[
  \F_{p^m}\cap\F_{p^n}
  =\F_{p^{\gcd(m,n)}}
\]
이고
\[
  \F_{p^m}\F_{p^n}
  =\F_{p^{\operatorname{lcm}(m,n)}}.
\]
여기서 왼쪽 두 번째 식은 두 체가 생성하는 합성체를 뜻한다.
\end{corollary}

\begin{proof}
교집합은 두 체의 공통 부분체이므로 그 차수는 $m$과 $n$의 공약수이다. 반대로 $d=\gcd(m,n)$이면 $d\mid m,n$이므로 $\F_{p^d}$는 두 체에 모두 들어간다. 부분체의 유일성으로 교집합은 $\F_{p^d}$이다.

합성체는 $\F_{p^m}$과 $\F_{p^n}$을 모두 포함하는 가장 작은 유한체이다. 그 차수는 $m$과 $n$의 공배수여야 하고, $\ell=\operatorname{lcm}(m,n)$이면 $\F_{p^\ell}$이 두 체를 모두 포함한다. 따라서 합성체는 $\F_{p^\ell}$이다.
\end{proof}

\begin{example}
$\overline{\F_p}$ 안에서
\[
  \F_{p^6}\cap\F_{p^{10}}=\F_{p^2},
  \qquad
  \F_{p^6}\F_{p^{10}}=\F_{p^{30}}.
\]
\end{example}

\subsection{유한체 위의 대수적 확장}

\begin{theorem}
유한체 $\F_q$의 모든 대수적 확장 $L/\F_q$는 다음 성질을 갖는다.
\begin{enumerate}[label=(\roman*)]
  \item $L/\F_q$는 분리확장이다.
  \item $f\in\F_q[X]$가 기약이고 $L$ 안에 한 근을 가지면, $f$는 $L[X]$에서 일차인수의 곱으로 완전히 분해된다.
\end{enumerate}
특히 모든 유한체는 완전체이다.
\end{theorem}

\begin{proof}
먼저 $E/\F_q$가 유한확장이고 $[E:\F_q]=n$이라 하자. 그러면 $|E|=q^n$이고 $E$는
\[
  X^{q^n}-X
\]
의 $\F_q$ 위 분해체이다. 이 다항식의 도함수는 $-1$이므로 분리다항식이다. 따라서 $E/\F_q$의 모든 원소는 $\F_q$ 위에서 분리적이다.

이제 $L/\F_q$가 임의의 대수적 확장이고 $\alpha\in L$라 하자. 유한체 $E=\F_q(\alpha)$의 원소 수를 $q^n$이라 쓰면, $\alpha$의 최소다항식은 분리다항식 $X^{q^n}-X$를 나누므로 분리된다. 따라서 $L/\F_q$는 분리확장이다.

또한 기약다항식 $f\in\F_q[X]$가 $L$ 안에 근 $\alpha$를 가지면 $f$는 $\alpha$의 최소다항식이다. 위에서 본 것처럼 $E=\F_q(\alpha)$는 $X^{q^n}-X$의 근 전체를 포함하며, $f$의 모든 근도 $E\subseteq L$에 속한다. 그러므로 $f$는 $L[X]$에서 완전히 분해된다. 모든 대수적 확장이 분리이므로 $\F_q$는 완전체이다.
\end{proof}

\begin{remark}
다음 장에서는 (ii)의 성질을 갖는 대수적 확장을 \emph{정규확장}이라 정의하고 여러 동치조건을 증명한다. 따라서 위 정리는 유한체의 모든 대수적 확장이 정규이면서 분리라는 사실을 미리 보여 준다.
\end{remark}

\subsection{유한체의 대수적 폐포}

유한체들을 중첩된 사슬로 만들기 위해
\[
  \F_p\subseteq\F_{p^{2!}}\subseteq\F_{p^{3!}}
  \subseteq\cdots
\]
를 생각하자. $n!\mid(n+1)!$이므로 각 포함이 성립한다.

\begin{theorem}
고정된 대수적 폐포 $\overline{\F_p}$ 안에서
\[
  \overline{\F_p}
  =\bigcup_{n\ge1}\F_{p^{n!}}.
\]
즉 오른쪽의 합집합은 $\F_p$의 대수적 폐포이다.
\end{theorem}

\begin{proof}
오른쪽을 $K$라 하자. 중첩된 체들의 합집합이므로 $K$는 체이고 $\F_p$ 위에서 대수적이다.

$f\in K[X]$가 비상수 다항식이라 하자. 계수는 유한개이므로 어떤 $m$에 대해 모두 $\F_{p^{m!}}$에 들어간다. $f$의 분해체는 이 유한체의 유한확장이므로 어떤 양의 정수 $r$에 대해 원소 수가
\[
  p^{m!r}
\]
인 유한체이다. 충분히 큰 $N$을 택하면 $m!r\mid N!$이고, 포함관계 정리에 의해 그 분해체는 $\F_{p^{N!}}\subseteq K$에 들어간다. 따라서 $f$는 $K$에서 완전히 분해된다. $K$는 대수적으로 닫혀 있으므로 $\F_p$의 대수적 폐포이다.
\end{proof}

\begin{keyidea}
고정된 표수 $p$에서 유한체의 세계는 양의 정수의 나눗셈 관계를 그대로 반영한다.
\[
  \F_{p^m}\subseteq\F_{p^n}
  \quad\Longleftrightarrow\quad m\mid n.
\]
부분체 문제는 약수 문제로, 교집합과 합성체 문제는 최대공약수와 최소공배수 문제로 바뀐다.
\end{keyidea}

\begin{checkpoint}
\begin{enumerate}[label=(\alph*)]
  \item $\F_{3^{18}}$의 부분체를 모두 적어라.
  \item $\F_{2^8}\cap\F_{2^{12}}$와 두 체의 합성체를 구하라.
  \item $\bigcup_{n\ge1}\F_{p^n}$을 고정된 대수적 폐포 안에서 생각하면 왜 체가 되는지 설명하라.
\end{enumerate}
\end{checkpoint}
